\chapter*{Введение}
\addcontentsline{toc}{chapter}{Введение}

Поведенческая разработка через тестирование (Behaviour-Driven Development, BDD) представляет собой современную методологию разработки программного обеспечения, объединяющую процессы проектирования, реализации и проверки качества программных решений. В условиях повышения сложности программных систем и необходимости эффективного взаимодействия между разработчиками, тестировщиками и заказчиками BDD выступает важным инструментом формализации требований посредством исполняемых спецификаций, описанных на естественном языке.

Задачи практической работы:
\begin{enumerate}
    \item изучить синтаксис языка Gherkin и его основные конструкции (Feature, Rule, Background, Scenario, Scenario Outline);
    \item разработать пользовательские истории (User Story) для функциональных возможностей приложения, реализующего поиск палиндромных подстрок;
    \item реализовать BDD-сценарии в декларативном и императивном стилях;
    \item реализовать Step Definitions на языке JavaScript с использованием фреймворка Cucumber;
    \item настроить инфраструктуру запуска BDD-тестов в JavaScript-проекте и генерацию отчётов о результатах тестирования.
\end{enumerate}

Практическая часть работы посвящена разработке набора BDD-тестов для приложения, реализующего алгоритм Манакера, предназначенный для эффективного поиска самой длинной палиндромной подстроки в строке за линейное время.